% TikZ 图：非匀速圆周运动的加速度分解
% Day 11：圆周运动
\begin{tikzpicture}[scale=0.86,>=Latex]
  \coordinate (O) at (-0.65,0);
  \coordinate (P) at ($(O)+(35:2.25)$);
  \coordinate (AtEnd) at ($(P)+(125:0.62)$);
  \coordinate (AnEnd) at ($(P)+(215:0.86)$);
  \coordinate (AEnd) at ($(AtEnd)+(215:0.86)$);

  \draw[thick] (O) circle (2.25);
  \filldraw[black] (O) circle (1.8pt);
  \node[below left] at (O) {$O$};
  \draw[dashed] (O) -- (P);
  \filldraw[red] (P) circle (3pt);
  \draw[densely dashed,gray] ($(P)+(-55:0.70)$) -- ($(P)+(125:1.12)$);

  % 速度沿质点所在处的圆周切线方向。
  \draw[->,blue,very thick] (P) -- ++(125:1.22)
    node[above left=2pt] {$\vect v$};

  % 加速度平行四边形：三个加速度都从圆上的质点位置出发。
  \draw[dashed,gray] (AtEnd) -- (AEnd);
  \draw[dashed,gray] (AnEnd) -- (AEnd);
  \draw[->,orange,thick] (P) -- (AtEnd)
    node[pos=0.50,right=3pt,fill=white,inner sep=1pt] {$\vect a_t$};
  \draw[->,red,very thick] (P) -- (AnEnd)
    node[pos=0.78,below right=2pt,fill=white,inner sep=1pt] {$\vect a_n$};
  \draw[->,purple,very thick] (P) -- (AEnd)
    node[pos=0.66,above left=2pt,fill=white,inner sep=1pt] {$\vect a$};
  \filldraw[red] (P) circle (3pt);

  \node[anchor=west,align=left,font=\footnotesize] at (3.20,0.30)
    {法向：$a_n=\dfrac{v^2}{r}$\\[4pt]
     切向：$a_t$ 改变速率\\[4pt]
     合加速度：$\vect a=\vect a_n+\vect a_t$};
\end{tikzpicture}
